\subsubsection*{Solution}
\paragraph*{Step-by-Step Derivation}

\subparagraph*{1. Action and background equations}

Using the signature $ds^2=dt^2-a^2(t)d\mathbf{x}^2$, the total first-order action is

\[
\mathcal S = \mathcal S_{\rm EH} + \mathcal S_{\vartheta} + \mathcal S_{\rm CS},
\]

where

\[
\mathcal S_{\rm EH} = -\frac{M_{\rm Pl}^2}{4} \int_{\mathcal M} \epsilon_{ABCD}\, e^A\wedge e^B\wedge R^{CD},
\]

\textcolor{blue}{and, for the pseudoscalar,}

\[
{\color{blue}
\mathcal S_{\vartheta} = \int_{\mathcal M} \left[ \frac12 \star d\vartheta\wedge d\vartheta \;-\; \frac{1}{4!}\,\epsilon_{ABCD}\, e^A\wedge e^B\wedge e^C\wedge e^D\; V(\vartheta) \right],}
\]

\textcolor{blue}{where $\star$ denotes the Hodge dual. Since $\tfrac{1}{4!}\epsilon_{ABCD}\,e^A\wedge e^B\wedge e^C\wedge e^D=\sqrt{-g}\,d^4x$ is the spacetime volume form, this is equivalently, in coordinate notation,}

\[
\mathcal S_{\vartheta} = \int d^4x\,\sqrt{-g} \left[ \frac12 g^{\mu\nu}\partial_\mu\vartheta\partial_\nu\vartheta -V(\vartheta) \right].
\]

The potential is interpreted exactly as printed in the problem,

\[
V(\vartheta)=\frac12m\vartheta^2,\qquad V_{,\vartheta}=m\vartheta,
\]

rather than as the frequently used convention $V=\frac12m^2\vartheta^2$.

These conventions and the resulting torsionful FLRW equations are given in the primary treatment of this model, \href{\detokenize{https://inspirehep.net/files/eb3973db1bbed05a53c1b06c37b78a17}}{Adshead, Brahma and Das, JCAP 03 (2026) 057}.

For

\[
e^0=dt,\qquad e^i=a(t)dx^i,
\]

the torsion ansatz gives the spin-connection components

\[
\bar\omega^{i0}=H e^i,\qquad \widetilde\omega^{0i}=h e^i,\qquad \widetilde\omega^{ij}=\phi\,\epsilon^{ij}{}_k e^k.
\]

Defining

\[
X\equiv H+h,
\]

the curvature components are

\[
R^{0i} = \left(\dot X+HX\right)e^0\wedge e^i + \phi X\,\epsilon^i{}_{jk}e^j\wedge e^k,
\]

\[
R^{ij} = \left(X^2-\phi^2\right)e^i\wedge e^j + (\dot\phi+H\phi)\epsilon^{ij}{}_k e^k\wedge e^0.
\]

Variation of the action then produces the background system [equations (3.12)-(3.16) in the cited paper]:

\[
X^2-\phi^2 = \frac{1}{3M_{\rm Pl}^2} \left(\frac12\dot\vartheta^2+V\right),
\]

\[
2(\dot X+HX)+X^2-\phi^2 = -\frac{1}{M_{\rm Pl}^2} \left(\frac12\dot\vartheta^2-V\right),
\]

\[
\phi = \frac{\alpha\dot\vartheta}{2M_{\rm Pl}^2} \left(X^2-\phi^2\right),
\]

\[
h = -\frac{\alpha\dot\vartheta\,\phi X}{M_{\rm Pl}^2},
\]

and

\[
\ddot\vartheta+3H\dot\vartheta+V_{,\vartheta} = -3\alpha \left[ \dot\phi(X^2-\phi^2) +2\phi X\dot X +3H\phi X^2-H\phi^3 \right].
\]

\subparagraph*{2. Algebraic elimination of torsion}

Set $M_{\rm Pl}=1$, $u=\dot\vartheta$, and

\[
E(\vartheta,u) \equiv \frac13\left(\frac12u^2+\frac12m\vartheta^2\right).
\]

The constraint equations give

\[
X^2-\phi^2=E,
\]

\[
\phi=\frac{\alpha u}{2}E.
\]

Consequently,

\[
X=\sqrt{E+\phi^2},
\]

where the positive root selects the expanding solution. Defining

\[
q\equiv\alpha u\phi =\frac{\alpha^2u^2E}{2},
\]

the remaining torsion equation gives

\[
h=-qX.
\]

Since $X=H+h$,

\[
H=X-h=(1+q)X.
\]

Thus the Hubble parameter used in the numerical integration is

\[
{ H(\vartheta,u) = \left(1+\frac{\alpha^2u^2E}{2}\right) \sqrt{ E+\frac{\alpha^2u^2E^2}{4} } }.
\]

\subparagraph*{3. Explicit scalar evolution equation}

Because $\phi$ and $X$ depend on $(\vartheta,u)$,

\[
\dot\phi=\phi_{,\vartheta}u+\phi_{,u}\dot u, \qquad \dot X=X_{,\vartheta}u+X_{,u}\dot u.
\]

The required derivatives are

\[
E_{,\vartheta}=\frac{m\vartheta}{3}, \qquad E_{,u}=\frac{u}{3},
\]

\[
\phi_{,\vartheta} = \frac{\alpha u}{2}E_{,\vartheta}, \qquad \phi_{,u} = \frac{\alpha}{2}\left(E+uE_{,u}\right),
\]

\[
X_{,z} = \frac{E_{,z}+2\phi\phi_{,z}}{2X}, \qquad z\in\{\vartheta,u\}.
\]

Writing

\[
B_0 = u\left(\phi_{,\vartheta}E+2\phi X X_{,\vartheta}\right) + H\phi(3X^2-\phi^2),
\]

\[
B_1 = \phi_{,u}E+2\phi X X_{,u},
\]

the Klein-Gordon equation becomes explicitly

\[
{ \dot u = \frac{-3Hu-m\vartheta-3\alpha B_0} {1+3\alpha B_1} }.
\]

The number of e-folds satisfies

\[
\dot N=H, \qquad N(0)=0,
\]

so that

\[
N(t)=\ln\frac{a(t)}{a(0)} =\int_0^t H(t')\,dt'.
\]

\subparagraph*{4. Initial conditions}

Using

\[
\alpha=10^{-4},\quad m=10^{-6},\quad \vartheta(0)=15,\quad u(0)=0.1,
\]

gives

\[
V(15)=\frac12(10^{-6})(15)^2 =1.125\times10^{-4},
\]

\[
E(0) = \frac{0.0051125}{3} = 0.00170416666666667,
\]

\[
\phi(0) = 8.52083333333333\times10^{-9},
\]

\[
H(0) = 0.0412815535883405.
\]

\subparagraph*{5. Numerical integration}

I integrated the system

\[
\dot\vartheta=u,\qquad \dot u=F(\vartheta,u),\qquad \dot N=H(\vartheta,u)
\]

with the adaptive eighth-order DOP853 method. Tightening the relative tolerance and reducing the maximum step verifies the quoted digits.

At $t=25000$, the solution is

\[
\vartheta=0.0486840706823,
\]

\[
\dot\vartheta = 2.25506937363\times10^{-4},
\]

\[
H = 9.41837897738\times10^{-5},
\]

and

\[
N(25000) = 75.1009281083.
\]

Thus the accumulated logarithmic scale-factor growth at the requested endpoint is $75.1009281083$ e-folds.
